\documentclass[10pt,a4paper]{amsart}
\usepackage[margin=19mm]{geometry}
\usepackage{amsmath,amssymb,mathtools}
\usepackage[scr=rsfs,cal=euler]{mathalfa}
\usepackage{xfrac}
\usepackage[unicode,hidelinks,bookmarks=false]{hyperref}
\newcommand{\ie}{i.e.\ }
\newcommand{\cf}{cf.\ }
\newcommand{\resp}{respectively\ }
\newcommand{\Subsection}{\subsection}
\providecommand{\nobreakdash}{\nobreak\hbox{-}}
\setlength{\emergencystretch}{2em}
\multlinegap0pt
\usepackage{bm}
\usepackage{bbm}
\usepackage{dsfont}
\usepackage{xspace}
\usepackage{cancel}
\usepackage{mathtools}
\usepackage{amsthm}
\makeatletter
\@ifundefined{thm}{\newtheorem{thm}{Theorem}}{}
\makeatother
\newcommand{\cP}{\mathcal{P}}
\renewcommand{\L}{\mathcal{L}}
\title[Topological line arrangements with high multiplicities]{Topological line arrangements with high multiplicities}
\author{Paolo Aceto, Marco Golla}
\date{Source: arXiv:2607.09398v1 (CC BY 4.0)}
\begin{document}
\maketitle

\noindent\emph{Source and licence.} Topological line arrangements with high multiplicities. Version arXiv:2607.09398v1 (CC BY 4.0), distributed under the Creative Commons Attribution 4.0 International licence, as recorded in the arXiv licence metadata for this exact version and shown on the abstract page https://arxiv.org/abs/2607.09398v1. Full rights notice: licenses/arxiv-2607.09398-CC-BY-4.0.txt.\par\smallskip

\noindent\textit{Editorial scope.} Counted unit: the Thm whose source label is \texttt{t:configurations} in the article (source0.tex lines 458--477, source label \texttt{t:configurations}), delivered together with the article's own complete original proof of it (source0.tex lines 492--571, one \texttt{proof} environment of 80 lines). No mathematics is written, repaired or re-derived: statement and proof are the article's own text line for line. The proof is detached from the statement in the source: the article states the result at lines 458--477 and prints its proof at lines 492--571, and the proof environment's own header names the statement, so the association is the article's own. Required context retained uncounted: t:divisible (lines 223--232). Every definition, display and label of those lines is retained unchanged. Omitted from this package and not redistributed: 1--222, 233--457, 478--491, 572--674. Nothing inside a retained excerpt is omitted or altered: every retained body is delivered exactly as the article's lines are written, so no omission receipt of any kind accompanies this package. Citations: the retained excerpts carry no citation command at all, so no bibliography entry of the article is transcribed and no reference is redistributed. Reference adaptation: none was needed and none was made; every reference inside the delivered unit resolves inside it, so no unresolved reference or citation is produced. Printed slips kept literal: no printed word, symbol, hypothesis or number of the counted statement or of its proof was altered. Layout and class: the article's own class is \texttt{amsart} with packages including inputenc, amsfonts, amsmath, amssymb, amsthm, babel, color, geometry, pinlabel; the wrapper is the lane's pinned \texttt{amsart} editorial wrapper at 10pt on A4, which restates the theorem-like environments the retained text needs and adds the article's own macro definitions that the retained text needs (\texttt{\textbackslash L}, \texttt{\textbackslash cP}), with extra wrapper packages (bm, bbm, dsfont, xspace, cancel, mathtools). The delivered numbering is the unit's own. Assets: no retained span contains a figure environment, a graphics-inclusion command, a PDF-page inclusion, a table or a TikZ picture; no image file is copied into this package, the package declares no asset dependency and no image file is redistributed with it. Licence: version arXiv:2607.09398v1 of this article is distributed under the Creative Commons Attribution 4.0 International licence, as recorded in the arXiv licence metadata for this exact version and shown on the abstract page https://arxiv.org/abs/2607.09398v1. A full rights notice is delivered at licenses/arxiv-2607.09398-CC-BY-4.0.txt. Characters: every retained excerpt is pure ASCII, so the delivered engine, XeLaTeX with its default Latin Modern text font, reports no missing character.
\begin{thm}\label{t:divisible}
Let $\L$ be a $q$-divisible line arrangement for some $q$. Then:
\[
\sum_m (m-{\textstyle \frac{6q}{q+1}})t_m \le d.
\]
In particular, if $q=2$ we have:
\[
\sum_{m \ge 6} (m-4)t_m \le d + 2t_2.
\]
\end{thm}

\begin{thm}\label{t:configurations}
If a line arrangement of degree $n$ and with $t_k \ge n$ is realisable, then:
\[
n \ge \left\{\begin{array}{ll}
k^2 - 4 & \mbox{if } n \equiv 0,\, k \equiv 0 \pmod 2\\
k^2 - 5 & \mbox{if } n \equiv 1,\, k \equiv 0 \pmod 2\\
k^2 - 5 & \mbox{if } n \equiv 0,\, k \equiv 1 \pmod 2\\
k^2 - 4 & \mbox{if } n \equiv 1,\, k \equiv 1 \pmod 2\\
\end{array}
\right.
\]
%\[
%d \ge k^2 - 4\frac{1+r_2(\cP,\L)}d > k^2-4,
%\]
Therefore, regardless of the parity of $n$ and $k$,
\[
n \ge k^2-5.
\]
In particular, this applies to $(n_k)$-configurations.
\end{thm}

\begin{proof}[Proof of Theorem~\ref{t:configurations}]
Let $(\cP, \L)$ be a realisable $(n_k)$-configuration.
As noted above, $n \ge k^2-k+1$, so the inequalities are already true combinatorially for $k < 6$.
So we can suppose $k \ge 6$.

By perturbing the points of multiplicities different than $k$ and $2$ if necessary, we obtain a topological realisation of a line arrangement $\L'$ with $t_k(\L') = n$ and $t_2(\L') = {n \choose 2} - n {k \choose 2}$.

By applying Theorem~\ref{t:divisible} to a $(d,t,\ell)$-arrangement, we obtain
\[
(\ell-4)t \le d + (d^2-d) - d(\ell^2-\ell),
\]
which is equivalent to
\begin{equation}\label{e:dtl}
d^2 \ge (\ell^2-4)t.
\end{equation}
We now construct a $(d,t,\ell)$-arrangement starting from $\L'$, in three of the four cases.

\begin{itemize}
\item If $n$ and $k$ are both even, the arrangement $\L'$ defined above is an $(n,n,k)$-arrangement, and~\eqref{e:dtl} simplifies to $n\ge k^2-4$, as required.

\item Suppose $n$ is odd and $k$ is even.
As on average a line in the arrangement contains $k$ points of multiplicity $k$, we can find a line in $\L'$ that contains at least $k$ points\footnote{Any line would work if we worked with a configuration.}.
Remove this line from the arrangement. 
All intersection points that were on this line either disappeared (because they were double points) or have now multiplicity $k-1$, which is odd. 
Perturb each of the latter points to obtain an $(n-1,t,t',k)$-arrangement with $t + t' = n$ and $t \le n-k$.
Applying Theorem~\ref{t:divisible} and $t \le n-k$ we obtain:
\[
(k-4)t + (k-6)t' \le n-1 + 2{d \choose 2} - 2t{k \choose 2} - 2t'{k-2 \choose 2}.
\]
Since $t \ge n-k$,
\[
(k-4)(n-k) + (k-6)k \le (k-4)t + (k-6)t'
\]
and
\[
2{d \choose 2} - 2t{k \choose 2} - 2t'{k-2 \choose 2} \le 2{n \choose 2} - 2(n-k){k \choose 2} - 2k{k-2 \choose 2},
\]
so that
\[
(k-4)(n-k) + (k-6)k \le n-1 + 2{n \choose 2} - 2(n-k){k \choose 2} - 2k{k-2 \choose 2}.
\]
From this one obtains:
\[
n^2 - (k^2-2)n + 4k^2-4k+1 \le 0,
\]
and solving the degree-2 inequality in $n$ yields
\[
n \ge \frac{k^2-2 + \sqrt{k^4-20k^2 + 16k}}2 > \frac{k^2-2 + k^2-12}2 = k^2-7,
\]
where we used $k \ge 6$ to say that $\sqrt{k^4-20k^2 + 16k} > k^2-12$.
Since $n$ is odd and $k$ is even, $n > k^2-7$ implies $n \ge k^2-5$.
(Note that the other half-line of solutions of the quadratic inequality is not relevant, since the other solution of the quadratic equation is smaller than $1$.)

\item If $n$ is even and $k$ is odd, we remove $k-1$ concurrent lines from $\L'$ and we perturb all remaining points of multiplicity $k$ into a point of multiplicity $k-1$ and $k-1$ points of multiplicity 2. The resulting arrangement is an $(n-k+1,n-1,k-1)$-arrangement, since we lowered the degree by $k-1$, we removed a single point of multiplicity $k$, and turned all others into points of multiplicity $k-1$ (either by removing a line through them or by perturbing). Equation~\eqref{e:dtl} now gives:
\[
(n-k+1)^2 \ge (k^2-2k-3)(n-1)  \Longleftrightarrow  n^2 - (k^2-5)n + 2k^2-4k-2.
\]
Solving for $n$, we get:
\[
n \ge \frac{k^2 - 5 + \sqrt{k^4-18k^2 + 16k + 33}}2 > \frac{k^2 - 5 + k^2- 9}2 = k^2-7.
\]
Here we used $k \ge 5$ to say that $16k+33 > 81$, which in turn implies $\sqrt{k^4-18k^2 + 16k + 33} > k^2-9$.
(Similarly to the previous case, the other solution of the quadratic equation is smaller than $4$, so not geometrically relevant.)
Since $n$ and $k$ have opposite parities, we actually obtain $n \ge k^2-5$.

\item If $n$ and $k$ are both odd, we remove all lines passing through a point of multiplicity $k$ and perturb all remaining points of multiplicity $k$ as above, to get an $(n-k,n-1,k-1)$-arrangement.
Equation~\eqref{e:dtl} yields:
\[
(n-k)^2 \ge (n-1)(k^2-2k-3) \Longleftrightarrow n^2 - (k^2-3)n + 2k^2-2k-3.
\]
We solve for $n$ and obtain:
\[
n \ge \frac{k^2 - 3 + \sqrt{k^4-14k^2 + 8k + 21}}2 > \frac{k^2 - 5 + k^2- 7}2 = k^2-6,
\]
where we used $k \ge 5$ to say that $8k+21 > 49$, and so $\sqrt{k^4-14k^2 + 8k + 21} > k^2-7$.
(As above, the other solution of the quadratic equation is smaller than 4, so not geometrically relevant.)
Since $n$ and $k$ have the same parity, we have $n \ge k^2-4$.
\end{itemize}
This concludes the proof.
\end{proof}

\end{document}
